\documentclass{article}
\usepackage{amsmath,amssymb,package-commands}

% Definitions stay in the source. The visual editor expands their math bodies.
\def\myparen#1{\left(#1\right)}
\newcommand{\R}{\mathbb{R}}
\newcommand{\norm}[1]{\left\lVert #1\right\rVert}
\newcommand{\pair}[2][x]{\langle #1,#2\rangle}
\newcommand{\diff}{\mathrm{d}}
\DeclareMathOperator{\rank}{rank}

\usepackage{wasysym,xcolor}
\definecolor{exampleblue}{RGB}{25,70,150}
\begin{document}
\section{Custom commands}

A plain TeX definition, \verb|\def\myparen#1{\left(#1\right)}|,
works with ordinary expressions and nested arguments:
$\myparen{x+1}$ and $\myparen{\frac{1}{\myparen{x}}}$.

The commands \verb|\R| and \verb|\norm{...}| combine standard math fonts
and paired delimiters: $x\in\R$, $\norm{x}$, and
$\norm{\frac{x}{y}}$.

The command \verb|\pair| has an optional first argument.
Its default is $\pair{y}$; an explicit argument gives $\pair[u]{v}$.

An operator declared with \verb|\DeclareMathOperator| has upright letters
and operator spacing: $\rank A=2$.

\[
\int_{\R} e^{-x^2}\,\diff x=\sqrt{\pi},
\qquad \mathcal{A}\subseteq\R.
\]

\begin{align*}
\norm{x+y}&\leq\norm{x}+\norm{y},\\
\myparen{a+b}^2&=a^2+2ab+b^2.
\end{align*}

\section{Editable definitions in the body}
Insert, Commands, Define command inserts an orange definition box.
Its initial name is selected. Rename it, then click its rendered body to edit.
These definitions are deliberately in the body so their fields can be tried:

\newcommand{\welcome}[2][Hello]{#1, \textbf{#2}!}

The default gives \welcome{reader}; an override gives \welcome[Welcome]{friend}.
Click the default field and leave it unchanged: the source uses the implicit
default. Change Welcome to Hello: its brackets should also disappear.
The source keeps an explicit argument only when it differs from the default.

\NewDocumentCommand{\orderedpair}{m O{y}}{\left(#1,#2\right)}

A required argument can precede an optional one: $\orderedpair{x}$ and
$\orderedpair{a}[b]$. Use Command definition, Configure definition arguments, or right-click an
indicator to change its type. Edit the default in its brackets or in the dialog.
Adding \verb|#3| to the body shows an inferred argument; remove it before
leaving to discard that provisional argument. Existing arguments require explicit deletion in the dialog.

\newcommand{\measure}[1]{The value is $\frac{#1}{2}$.}

Mixed text and mathematics stay editable in the definition: \measure{x+1}
Use the source toggle for literal changes and return to instant mode afterward.
For a command invocation anywhere in body text, Home, Editing, Edit command
and arguments opens its name and ordered argument rows.
\section{Commands discovered from packages}
\LaTeX{} loads the definitions in the adjacent package file.
PasTeX also reads these definitions for the instant preview:
$\packageparen{x+1}$, $\packagealias{y-1}$, and $\packagehalf{a+b}$.

Installed packages provide font declarations as well as macros:
\[
a\shortleftarrow b,\qquad x\boxdot y,\qquad \packageop(A,B).
\]
The arrow and boxed operator use the font declared by the installed
\texttt{stmaryrd} package.

\renewcommand{\packagehalf}[1]{\frac{#1}{3}}
This body definition overrides the package definition for the formulas below:
$\packagehalf{x+y}$.

\def\bodyparen#1{\left[#1\right]}
Plain \TeX{} definitions can also be edited here: $\bodyparen{x+1}$.

\section{Groups and conditional text}
\begin{abstract}
An abstract is editable native text, including formulas such as $E=mc^2$.
\end{abstract}
\begingroup\color{exampleblue}This color is local to the explicit group.\endgroup
This text uses the surrounding color. The \verb|\Square| symbol is \Square.

\ifcsname myparen\endcsname
The earlier command is defined.
\else
The command is undefined.
\fi
\if aa The character tokens match.\else They differ.\fi
\end{document}
